% !TeX program = xelatex
% 文档：SmarTAI 题目识别与判卷问题记录（基于 MATH—OCR/test.pdf 测试）
\documentclass[12pt,a4paper]{ctexart}

\usepackage[margin=2.5cm]{geometry}
\usepackage{booktabs}
\usepackage{enumitem}
\usepackage{amsmath}

\title{\bfseries SmarTAI 题目识别与判卷问题记录}
\author{测试来源：\texttt{SmarTAI-shared/MATH--OCR/test.pdf}}
\date{\today}

\begin{document}
\maketitle

\section*{问题概述}

以下问题均在以 \texttt{test.pdf} 为题目源、进行题目识别（question preparation）与
判卷（grading）时观察到。问题按严重程度排序。

\section{选择题、填空题不应有过程分}

\subsection{现象}
当前判卷逻辑对选择、填空类题目评分时仍可能产生“过程分”。这类题目答案唯一，
评分应只依据最终结果是否正确，不应存在中间步骤得分。

\subsection{期望}
\begin{itemize}[leftmargin=2em]
  \item 选择题、填空题只判结果分：正确得满分，错误得 0 分；不输出过程分项。
  \item 若题干要求写出简要过程，则需在题目类型与评分规则中显式标注，
        默认不启用过程分。
\end{itemize}

\section{重复识别问题：同一道题被识别为两个题目}

\subsection{现象}
对 \texttt{test.pdf} 的识别结果存在两类重复：

\begin{enumerate}[leftmargin=2em,label=(\arabic*)]
  \item \textbf{解答题中的“第一问”被单独列出为一个题目}，
        同时完整整道题又被识别为一个题目，造成同一题两份；
  \item \textbf{选择、填空题也出现同一道题被重复识别}的情况。
\end{enumerate}

\subsection{影响}
\begin{itemize}[leftmargin=2em]
  \item 题目总数虚高（如本该 19 题却出现22条目）；
  \item 同一题设立两个 \texttt{q\_id}，学生对同一题分别作答两份，判卷结果重复、混乱；
  \item 展示与提交阶段无法区分“小题”与“整题”。
\end{itemize}

\subsection{期望}
\begin{itemize}[leftmargin=2em]
  \item 解答题整体作为一个题目，子问（(1)(2)(3)）作为其内部结构保留，
        不作为独立题目条目；
  \item 去重：同一题号、同一题面只保留一个题目；
  \item 若需要小题级评分，应通过题目内部的小题结构表达，而非拆成独立题目。
\end{itemize}

\section{满分分配混乱（由问题 2 引发）}

\subsection{现象}
由于同一道题（整题与第一问）被重复识别，其 \texttt{max\_score}
（满分）也被重复设置；两个条目各自带一份满分，导致：
\begin{itemize}[leftmargin=2em]
  \item 单题满分与整套卷面总分不匹配；
  \item 判卷时同一题被评两次，分数被重复计入或相互覆盖；
  \item 教师端“各题总分不同”的配置项与识别结果不一致（见
        \texttt{问题1总分 (各题总分不同).png}）。
\end{itemize}

\subsection{期望}
\begin{itemize}[leftmargin=2em]
  \item 每道题只存在一个 \texttt{max\_score}，由教师统一配置；
  \item 识别阶段不擅自为重复条目生成满分；去重后满分只落在唯一题目上；
  \item 总分统计应基于去重后的题目集合。
\end{itemize}

\section{判卷过程大量超时}

\subsection{现象}
判卷阶段频繁出现“所有专家均网络/超时错误，暂未批改”，表现为：
\begin{itemize}[leftmargin=2em]
  \item \texttt{error\_code = provider\_timeout}（部分环境为
        \texttt{provider\_unreachable}）；
  \item 一次判卷中多个题目/多个专家同时报超时；
  \item 重试后偶可成功，但高峰期大面积失败。
\end{itemize}

\subsection{初步定位}
\begin{itemize}[leftmargin=2em]
  \item 当前模型经 USTC 免费中转（\texttt{api.llm.ustc.edu.cn}）调用，
        该中转对较大请求可能挂起（此前识别阶段单次大请求可挂起约 600 秒）；
  \item 判卷并发扇出多道题、多专家，放大网络波动的影响；
  \item 与提示词内容基本无关（同一提示词在可用时段可成功判卷）。
\end{itemize}

\subsection{期望}
\begin{itemize}[leftmargin=2em]
  \item 缩短单次模型调用超时（自定义端点超时），使失败快速进入重试/报错路径；
  \item 对识别阶段大输入做分块（已实现于本地分支），控制单请求体积；
  \item 判卷侧增加失败重试/退避，避免一次超时即整批失败；
  \item 长期建议接入官方 API（DeepSeek/Qwen 官方域名），绕开中转不稳定。
\end{itemize}

\section*{复现步骤（供回归）}
\begin{enumerate}[leftmargin=2em]
  \item 上传 \texttt{test.pdf} 为题目源；
  \item 触发题目识别（question preparation），检查：题目总数、是否重复、
        子问是否被拆成独立题、各题满分；
  \item 提交学生作答并判卷，观察：选择/填空是否出现过程分、判卷是否超时。
\end{enumerate}

\section*{结论}
问题 2（重复识别）为根因，直接导致问题 3（满分混乱）；
问题 1（过程分）为评分规则缺失；问题 4（超时）为模型中转稳定性问题。
建议按顺序处理：先修复去重与子题结构，再规范评分规则，最后优化判卷网络稳健性。

\end{document}
